\chapter{The atomic model}
\label{ch:atomic}

\section{Electronic states}

The electronic state of an atom is its \emph{shell configuration}: three integers \code{K L M}, the number of
electrons in the $n=1,2,3$ shells (capacities 2, 8 and 18).
There are no subshells and no terms or levels within a configuration; whatever an atomic-structure code resolves
is summed or averaged per configuration when the data are generated (Chapter~\ref{ch:data}).
An atom starts in its element's ground configuration (table~\ref{tab:elements}).
Its charge in the MD is $q = Z - (K+L+M)$: at the first step the module replaces the force field's partial charges by
these integer ionization charges (all zero at the start), and every transition changes the configuration and the
charge.

An atom's element is the one whose standard atomic weight is nearest to its mass, within 0.5 (so argon, 39.948, and
calcium, 40.078, are told apart):

\begin{table}[H]
\centering\small
\begin{tabular}{@{}lrrl@{\quad}lrrl@{\quad}lrrl@{}}
\toprule
element & mass & $Z$ & ground & element & mass & $Z$ & ground & element & mass & $Z$ & ground \\
\midrule
H  & 1.008  & 1 & \code{1 0 0} & F  & 18.998 & 9  & \code{2 7 0} & S  & 32.06  & 16 & \code{2 8 6} \\
Li & 6.94   & 3 & \code{2 1 0} & Ne & 20.180 & 10 & \code{2 8 0} & Cl & 35.45  & 17 & \code{2 8 7} \\
Be & 9.0122 & 4 & \code{2 2 0} & Na & 22.990 & 11 & \code{2 8 1} & Ar & 39.948 & 18 & \code{2 8 8} \\
B  & 10.81  & 5 & \code{2 3 0} & Mg & 24.305 & 12 & \code{2 8 2} & Ca & 40.078 & 20 & \code{2 8 10}$^\ast$ \\
C  & 12.011 & 6 & \code{2 4 0} & Al & 26.982 & 13 & \code{2 8 3} & Fe & 55.845 & 26 & \code{2 8 16}$^\ast$ \\
N  & 14.007 & 7 & \code{2 5 0} & Si & 28.085 & 14 & \code{2 8 4} & Ni & 58.693 & 28 & \code{2 8 18}$^\ast$ \\
O  & 15.999 & 8 & \code{2 6 0} & P  & 30.974 & 15 & \code{2 8 5} & I  & 126.904 & 53 & \code{2 8 18}$^{\ast\ast}$ \\
\bottomrule
\end{tabular}
\caption{Elements the module knows (\code{mcmd\_elements[]} in \file{atomic\_data.c}); H to Ar have built-in data.
$^\ast$\,The 4s electrons are counted in M; only with \opt{mcmd-atomic-data = files}.
$^{\ast\ast}$\,A placeholder: right for the force field, not for the ionization physics.
Deuterium (mass 2) is hydrogen.  Helium is left out: an atom of mass $\sim$4 is far more likely a heavy hydrogen
(\code{pdb2gmx -heavyh}, mass repartitioning) and is refused rather than taken as helium.}
\label{tab:elements}
\end{table}
A new element needs a row in that table (mass, symbol, ground configuration) and, above argon or for files, its
atomic data.

\section{Transitions}

Three processes are tabulated per configuration, each with its final configuration:
\begin{description}
  \item[Photoionization] (type 2) removes one electron from a shell.
    The table gives the cross section $\sigma$ (cm$^2$) at the photon energy; the rate is $\sigma I(t)$ with the pulse
    flux $I(t)$ below.
  \item[Auger--Meitner decay] (type 0) fills a hole from an outer shell and ejects a second electron: the atom loses
    one electron.
    The table gives the rate (s$^{-1}$).
  \item[Fluorescence] (type 1) fills a hole from an outer shell and emits a photon, which leaves; the charge does not
    change.
\end{description}
Charge transfer (Chapter~\ref{ch:ct}) and, with the electron cloud, collisional ionization
(Chapter~\ref{ch:coll}) are added to these.

\section{The pulse}

The pulse is Gaussian in time, with peak time $t_0$ (\opt{mcmd-pulse-peak-time}), full width at half maximum
$\tau$ (\opt{mcmd-pulse-fwhm}), $N_{\mathrm{ph}}$ photons (\opt{mcmd-pulse-photons}) in a focal spot of
diameter $d$ (\opt{mcmd-pulse-focal-diameter}), uniform over the sample:
\begin{equation}
  I(t) = \frac{N_{\mathrm{ph}}}{\pi (d/2)^2}\,\frac{1}{\sigma_t\sqrt{2\pi}}\,
         \exp\!\Big[-\frac{(t-t_0)^2}{2\sigma_t^2}\Big],\qquad \sigma_t = \frac{\tau}{2\sqrt{2\ln 2}},
\end{equation}
in photons cm$^{-2}$\,s$^{-1}$.
The fluence $N_{\mathrm{ph}}/(\pi d^2/4)$ is what sets the ionization; the example's $10^{11}$ photons into 100\,nm
are $1.27\times10^{13}$ photons/\textmu m$^2$.
The photon energy (\opt{mcmd-pulse-photon-energy}) only enters through the atomic data, which are computed for it
(built in) or must be generated for it (files),
and through the energy of the photoelectrons.
\opt{mcmd-pulse-bandwidth} (relative FWHM, default 0.003) gives each photoelectron a Gaussian jitter of that width
in its energy, which matters only for the escape test near threshold; the cross sections are those of the nominal
energy.

\section{The Monte Carlo step}

Every MD step $\Delta t$, every atom is advanced by a first-reaction kinetic Monte Carlo~\cite{gillespie1976}:
for each transition available from its current configuration a waiting time $\delta t = -\ln u / k$ is drawn
($u$ uniform, $k$ the rate, with $I(t)$ for photoionization); the earliest fires if it falls within the time left in
the step; then the atom's new configuration offers new transitions, and the procedure repeats until the step is used
up.
Several transitions can therefore happen to one atom in one step, e.g.\ a K-shell photoionization followed within
femtoseconds by an Auger decay.
The Monte Carlo runs on the master rank from one random stream (seeded by \opt{mcmd-seed},
Section~\ref{sec:seed}), so the sequence of events does not depend on the number of ranks.

\section{The released electrons}

Each photoionization and Auger decay releases an electron, with kinetic energy from the configurations' total
binding energies $T$ (Chapter~\ref{ch:data}):
\begin{equation}
  \varepsilon_{\mathrm{photo}} = E_{\mathrm{ph}} - \big(T_{\mathrm{from}} - T_{\mathrm{to}}\big),\qquad
  \varepsilon_{\mathrm{Auger}} = T_{\mathrm{to}} - T_{\mathrm{from}}.
\end{equation}
For carbon at 1000\,eV a K-shell photoelectron has $\sim$700\,eV and a KLL Auger electron $\sim$260\,eV.
Without the electron cloud (\opt{mcmd-electrons = none}) every electron leaves.
With it, the electron escapes if $\varepsilon$ exceeds the depth of the potential well at its atom, and is trapped in
the cloud otherwise (Section~\ref{sec:escape}).
